\section{Evaluation}

We evaluate \textsc{Simulon} along two axes. First, we validate each
simulator component independently against controlled measurements on a
physical testbed, so that errors cannot cancel across subsystems.
Second, to assess whether the simulator generalises to configurations
and scales beyond the testbed, we compare end-to-end throughput
predictions against publicly reported numbers from the NVIDIA
Megatron-Bridge benchmark suite~\cite{megatron-bridge-perf}, which
covers a range of production LLM configurations on up to
$[\textbf{N\_MAX}]$ GPUs.

\subsection{Testbeds}

\paragraph{Physical testbed.}
Component-level validation is performed on a cluster of
$[\textbf{NUM\_NODES}]$ nodes, each equipped with
$[\textbf{NUM\_GPUS\_PER\_NODE}]$ NVIDIA $[\textbf{GPU\_MODEL}]$
GPUs ($[\textbf{GPU\_MEM}]$\,GB HBM, $[\textbf{GPU\_TF\_BF16}]$
TFLOP/s BF16). Within each node, GPUs are connected via NVSwitch with
an all-to-all bisection bandwidth of
$[\textbf{INTRA\_BW}]$\,GB/s per GPU. Nodes are
interconnected through a $[\textbf{SWITCH\_MODEL}]$ switch fabric
arranged in a $[\textbf{TOPOLOGY}]$ topology
(e.g.\ fat-tree with $[\textbf{NUM\_RAILS}]$ rails), providing
$[\textbf{INTER\_BW}]$\,GB/s inter-node bandwidth per GPU via
$[\textbf{NIC\_MODEL}]$ NICs. The collective library used throughout
is NCCL $[\textbf{NCCL\_VERSION}]$.

\paragraph{Megatron-Bridge benchmarks.}
For scale validation we use the NVIDIA Megatron-Bridge performance
summary~\cite{megatron-bridge-perf}, which reports
tokens\,s$^{-1}$\,GPU$^{-1}$ and model TFLOP\,s$^{-1}$\,GPU$^{-1}$
for a suite of production models---including LLaMA-3 (8B--405B),
DeepSeek-V3, and GPT-class models up to 120B parameters---across
DGX~H100, DGX~B200, and DGX~GB200 systems at scales ranging from 8
to $[\textbf{N\_MAX\_BRIDGE}]$ GPUs. Because these results are
third-party and cover hardware we do not operate directly, they
serve as an external check on end-to-end simulator accuracy rather
than a controlled component benchmark; we discuss the additional
sources of uncertainty this introduces in
Section~\ref{sec:e2e-validation}.

\paragraph{Validation methodology.}
Our validation strategy is designed to localise errors: if the three
component models are validated independently, a discrepancy in an
end-to-end comparison can be attributed to a specific subsystem rather
than absorbed by compensating errors. The three components and their
validation paths are:

\begin{enumerate}
  \item \textbf{Compute model} (Section~\ref{sec:compute-validation}).
        We run a single-GPU Megatron-LM training step (PP\,=\,TP\,=\,DP\,=\,1)
        and record per-phase wall times (\texttt{forward-compute},
        \texttt{backward-compute}, \texttt{optimizer}) at
        \texttt{timing\_log\_level=1}. With no parallelism, the
        simulator's prediction reduces to a sum of profiled kernel
        durations with no collective or scheduling logic, so any
        residual error is unambiguously in the compute model.

  \item \textbf{Collective model} (Section~\ref{sec:comm-validation}).
        We run \texttt{nccl-tests} for AllReduce, AllGather,
        ReduceScatter, and AllToAll at the exact message sizes that
        \textsc{Simulon} generates for a given configuration, and
        compare measured bandwidth and latency against the simulator's
        collective timing predictions. This is fully decoupled from
        compute.

  \item \textbf{Scheduling model} (Section~\ref{sec:pipeline-validation}).
        With PP\,>\,1 the 1F1B schedule yields a theoretically known
        pipeline-bubble ratio:
        \[
          \rho = \frac{p - 1}{p + m - 1},
        \]
        where $p$ is the number of pipeline stages and $m$ the
        number of microbatches. We measure the actual GPU idle fraction
        from Megatron's per-stage timers and compare it against the
        replayer's output. A match here confirms the scheduling logic
        independently of whether compute or communication timings are
        exact.
\end{enumerate}

Finally, Section~\ref{sec:e2e-validation} compares end-to-end
throughput predictions against the Megatron-Bridge benchmarks to
assess generalisation beyond the testbed.
